%% ch03 --- Feedback: why nothing grows forever
%%
%% Stub, generated by tools/gen_stubs.py from tools/chapters.json.
%% The brief below is the contract for this chapter: what it must argue, what
%% it must buy the reader, and what it must NOT take from its neighbours.
%% Writing the chapter means deleting the stub block below (the one that opens
%% on the next \chapter line) and replacing it with the chapter text -- after
%% which this file is never regenerated.
%% If the brief turns out to be wrong once the code has been run, change it in
%% the manifest and record the contradiction in CLAUDE.md under *Resolved
%% questions*.

\chapter{Feedback: why nothing grows forever}
\label{ch:feedback}

\chapterstub{%
Argument: real growth slows because the thing that grows uses up what it
grows on \dash{} feedback \dash{} and feedback makes a rule NONLINEAR: the
effect of a change depends on where you already are. Verhulst's logistic
growth (1838) and its S-curve, carrying capacity, and the difference between
a straight-line rule and a curved one: for a linear rule two causes add
(superposition), for a nonlinear one they do not. Local stability of a fixed
point on a curve: zoom in until the curve looks straight, and the slope
there \dash{} measured numerically, as how much a tiny nudge is multiplied
by one step \dash{} decides whether nudges grow or die. This is the slope
idea the rest of the book leans on, taught as a measurement rather than as
calculus. The chapter ends on a discrete logistic population that settles
smoothly for a gentle growth rate and then, turned up, overshoots and
oscillates \dash{} the hook for Chapter 5. Payoff: the reader can find the
fixed points of a nonlinear rule, estimate the slope at each by computing,
and predict which one the system approaches. Traps to elicit: the effect of
two pushes is the sum of the effects (only for linear rules); a limit to
growth means growth stops smoothly (it can overshoot and oscillate).
Listings: continuous logistic growth stepped finely; numerical slope at a
fixed point; the discrete logistic at two growth rates. Labs: estimate a
slope numerically and classify a fixed point. Figures: exponential against
logistic growth; the overshoot.

\medskip
\textbf{Word budget:} 3500.
}
